\chapter{Preamble}


\Peano\ is an open source framework for solvers on dynamically adaptive
Cartesian meshes.
Its core is built with C++, but many tools around it are written in Python.
\Peano\  is based upon the fact that spacetrees, a generalisation of the
classical octree concept, yield a cascade of adaptive Cartesian grids.
Consequently, any spacetree traversal is equivalent to an element-wise
traversal of the hierarchy of the adaptive Cartesian grids. The software
\Peano\  realises such a grid traversal and storage algorithm, and it provides
hook-in points for applications performing per-element, per-vertex, and so
forth operations on the grid. It also provides interfaces for dynamic load
balancing, sophisticated geometry representations, and other features.


Peano is currently available in its fourth generation. 
The development of the original set of Peano codes started around 2002.
2005-2009, we merged these codes into one Peano kernel (2nd generation). 
In 2009, I started a complete reimplementation of the kernel with special
emphasis on reusability, application-independent design and the support for rapid prototyping. 
This third generation of the code ended around 2019 when we released the
ExaHyPE code---a hyperbolic equation system solver engine which uses Peano's
AMR meshes.
The current guidebook discusses \Peano\ , which reuses ideas and lessons learned
from Peano~3 as well as lots of code buildling blocks, but can be seen as a
major rewrite starting from scratch.



\Peano\ has been the code base for multiple projects over the past years.
The most prominent one is likely ExaHyPE.
The first generation of ExaHyPE has been a stand-alone projects using Peano
as one required library or framework out of many.
The second generation, \ExaHyPE, is merged into \Peano as spin-off.
This document therefore covers \ExaHyPE, too.




\section*{Dependencies and prerequisites}

\Peano's core is plain C++17 code. 
We however use a whole set of tools around it:

\begin{itemize}
  \item GNU autotools (automake) to set up the system (required).
  \item C++17-compatible C++ compiler (required).
  \item Python3 (optional; not required if you work only with the C++
  baseline).
  \item MPI2 (optional). MPI's multithreaded support is required.
  \item Intel's Threading Building Blocks (TBB) or OpenMP 5. We also can work with C++
  threads, but found their performance suboptimal. TBB is not shipped as part of OneAPI.
  \item The Visualization Toolkit (VTK) if you want to use the built-in
  visualisation facilities (optional).
  \item Doxygen if you want to create HTML pages of PDFs of the in-code
  documentation.
\end{itemize}

\begin{remark}
I test and maintain \Peano\  for Linux only.
If you prefer Windows or MacOS, it should work as long as you provide the
mandatory tools from above, but I won't be able to help.
\end{remark}


\begin{center}
 \begin{tabular}{|l|l|p{8cm}|}
  \hline
   Software & Version & Remarks \\
  \hline
   GNU g++ & 7.5 or older & Works with \Peano\ core but not with some extensions
   such as \ExaHyPE. \\
   GNU g++ & 9.3.1 & Current development version. Works. \\
  \hline
   Intel icpc & 19.1 & Works with \Peano\ and all
   extensions such as \ExaHyPE. Tests with Intel Fortran have been successful,
   too. \\
  \hline
   VTK & 7--9 & Supported by \Peano. Default is 8, i.e.~if you need other
   version, you have to inform configure (see Chapter \ref{chapter:postprocessing}). \\
   \hline
   Python & 3 & Development version. Python 2 is not supported by frontend. \\
   \hline
 \end{tabular}
\end{center}


\section*{How the text is organised}

This \LaTeX~documentation first summarises some installation options.
We then present some high-level tutorials for some applications built on top of
\Peano.
This notably covers \ExaHyPE.
We continue with a discussion of \Peano's Python API, before I discuss C++
options and how to extend the integration of Python with the C++ framework.


While the tutorial present a user's view, i.e.~how to use pre-manufactured
ingredients, the latter chapters discuss how to implement new ingredients.
However, they still give a high-level overview and focus on rationale and design
decisions.
\Peano's core philosophy is \emph{documentation is in the code} (see Section
\ref{section:installation:documentation}).


In line with this philosophy, also the \LaTeX\ documentation is held within the
repository.
If you feel something is urgently missing and deserves a chapter or
things remain unclear, please extend the \LaTeX~documentation and issue a merge
request.


{
  \flushright
  \today 
  \\ 
  Tobias Weinzierl 
  \\
}

 
